\documentclass[../readings.tex]{subfiles}

\begin{document}

\subsection{Linear Operators and the Poisson Equation}
\label{sub:operators}

\textBox{%
Differential equations become matrix problems after discretization. The finite matrix should preserve the main structure of the operator: sparsity, symmetry, definiteness, spectrum, and conditioning.
}

\subsubsection{Differential Operators and Discretization}

\addDef{Linear Operator}{%
Map $\mathcal{A}: X \to Y$ between function spaces satisfying:
\begin{align*}
    \mathcal{A}(\alpha u + v) = \alpha \mathcal{A}u + \mathcal{A}v.
\end{align*}
Matrices are finite-dimensional operators; $\frac{d}{dx}$ and $\nabla^2$ are infinite-dimensional.
}
\label{def:linear-operator}

\addDef{Dirichlet Poisson Problem}{%
The one-dimensional Poisson problem with homogeneous Dirichlet boundary data is
\begin{align*}
    -u''(x)=f(x),\qquad 0<x<1,\qquad u(0)=u(1)=0.
\end{align*}
The boundary conditions are part of the problem data.
}
\label{def:dirichlet-poisson}

\addDef{Grid Function}{%
For $n$ interior grid points, let $h=1/(n+1)$ and $x_i=ih$, $i=1,\dots,n$. A grid function is the vector
\begin{align*}
    \bu=(u_1,\dots,u_n)^T,\qquad u_i\approx u(x_i).
\end{align*}
The boundary values are fixed separately by the boundary conditions.
}
\label{def:grid-function}

\addDef{Finite Difference Stencil}{%
Approximating $u''(x)$ on a grid with spacing $h$:
\begin{align*}
    u''(x_k) \approx \frac{u(x_{k-1}) - 2u(x_k) + u(x_{k+1})}{h^2}.
\end{align*}
\textbf{Accuracy:}\enspace Second-order $O(h^2)$ via Taylor expansion.
}
\label{def:finite-difference-stencil}

\addExcr{%
\begin{enumerate}
    \item Taylor expand $u(x+h)$ and $u(x-h)$ about $x$.
    \item Add the two expansions.
    \item Solve for $u''(x)$.
    \item Identify the leading error term and prove the $O(h^2)$ claim in Def~\ref{def:finite-difference-stencil}.
\end{enumerate}
}
\label{ex:finite-difference-stencil}

\addDef{1-D Discrete Laplacian}{%
Discretizing $-d^2/dx^2$ on $[0,1]$ with $u(0)=u(1)=0$ yields the tridiagonal system $\bT \mathbf{u} = \mathbf{f}$:
\begin{align*}
    \bT = \frac{1}{h^2} \text{tridiag}(-1, 2, -1) \in \fR^{n \times n}.
\end{align*}
}
\label{def:discrete-laplacian-1d}

\addRemark{%
\textbf{(Sparsity)}\enspace Each row of $\bT$ uses only the grid point and its two neighbors. This locality is why finite difference matrices are sparse.
}

\addThm{Spectrum Convergence}{%
Eigenvalues of $\bT$ are $\lambda_k = \frac{4}{h^2} \sin^2(\frac{k\pi}{2(n+1)})$.
As $n \to \infty$, $\lambda_k \to k^2 \pi^2$ (eigenvalues of continuous $-d^2/dx^2$).
}
\label{thm:laplacian-eigenvalues}

\addExcr{%
\begin{enumerate}
    \item Let $(\bv_k)_j=\sin(jk\pi h)$ with $h=1/(n+1)$.
    \item Compute the $j$-th entry of $\bT\bv_k$.
    \item Use $\sin(A-B)+\sin(A+B)=2\sin A\cos B$.
    \item Use $1-\cos\theta=2\sin^2(\theta/2)$.
    \item Derive the eigenvalue formula in Thm~\ref{thm:laplacian-eigenvalues}.
    \item Let $h\to 0$ with fixed $k$ and recover $k^2\pi^2$.
\end{enumerate}
}
\label{ex:laplacian-eigenvalues}

\addThm{Discrete Poisson Structure}{%
The matrix $\bT$ in Def~\ref{def:discrete-laplacian-1d} is sparse, symmetric positive definite, and has condition number
\begin{align*}
    \kappa_2(\bT)=O(h^{-2}).
\end{align*}
}
\label{thm:poisson-structure}

\addPrf{%
Sparsity and symmetry are immediate from the tridiagonal stencil. By Thm~\ref{thm:laplacian-eigenvalues}, all eigenvalues are positive because $1\leq k\leq n$ and $0<k\pi/(2(n+1))<\pi/2$. Thus $\bT$ is SPD. Also
\begin{align*}
    \lambda_{\min}\sim \pi^2,
    \qquad
    \lambda_{\max}\sim \frac{4}{h^2},
\end{align*}
so $\kappa_2(\bT)=\lambda_{\max}/\lambda_{\min}=O(h^{-2})$.
}

\addExcr{%
\begin{enumerate}
    \item Use the stencil to count the number of nonzeros in $\bT$.
    \item Use the quadratic form $\bu^T\bT\bu$ to show positive definiteness for nonzero interior grid functions.
    \item Use Thm~\ref{thm:laplacian-eigenvalues} to estimate $\lambda_{\min}$ and $\lambda_{\max}$.
    \item Derive the condition number estimate in Thm~\ref{thm:poisson-structure}.
\end{enumerate}
}
\label{ex:poisson-structure}

\subsubsection{Operator Norms and Conditioning}

\addDef{Boundedness}{%
An operator is \textbf{bounded} if $\|\mathcal{A}u\|_Y \leq C \|u\|_X$ for some $C < \infty$.
\begin{itemize}
    \item Matrices are always bounded.
    \item Differential operators are \textbf{unbounded}.
\end{itemize}
}

\addRemark{%
\textbf{Numerical Consequence:}\enspace Because the continuous Laplacian is unbounded, its discretization $\bT$ becomes increasingly ill-conditioned as $h \to 0$:
\begin{align*}
    \kappa(\bT) = \frac{\lambda_{\max}}{\lambda_{\min}} \approx \frac{4/h^2}{\pi^2} = O(h^{-2}).
\end{align*}
Refining the grid $2\times$ quadruples the condition number.
}

\addRemark{%
\textbf{(Solver consequence)}\enspace Dense Gaussian elimination ignores sparsity. For Poisson matrices, sparse Cholesky, CG with preconditioning, multigrid, or FFT-based solvers are the natural algorithms, depending on geometry and boundary conditions.
}

\subsubsection{Exercises}

\addExcr{%
\begin{enumerate}
    \item Build $\bT$ for $n=50$ and plot its eigenvalues vs.\ $k^2 \pi^2$ using Thm~\ref{thm:laplacian-eigenvalues}.
    \item Show that $\kappa(\bT) \to \infty$ as $n \to \infty$ using Thm~\ref{thm:poisson-structure}.
    \item Implement a 2-D Poisson solver on a unit square. Verify $O(N^{3/2})$ cost for direct sparse solve.
    \item Prove the differentiation operator is unbounded on $L^2[0,1]$ using $u_k(x) = \sin(k\pi x)$.
\end{enumerate}
}

\end{document}
